%% =============================================================================
%% flowshop.tex
%% Programación de líneas de producción: Flow shop permutacional (F||Cmax)
%% Formato APA 7.ª edición (bibliografía manual con \begin{thebibliography}-like)
%% Compilar con:  pdflatex flowshop.tex  (dos veces)
%%
%% FIGURAS (carpeta fig/ junto a este archivo):
%%   - Las figuras 3, 5, 6 y 7 se generan con:  python3 generar_figuras.py
%%   - Las figuras 1 y 4 están dibujadas en TikZ dentro de este archivo.
%%   - La figura 2 (fotografía real) debe descargarse de Wikimedia Commons;
%%     ver las instrucciones en el recuadro [FIGURA PENDIENTE] que aparece en el PDF.
%% Si falta algún archivo de imagen, el documento compila igual y muestra un
%% recuadro indicando de dónde obtenerlo.
%% =============================================================================
\documentclass[12pt,letterpaper]{article}

\usepackage[T1]{fontenc}
\usepackage[utf8]{inputenc}
\usepackage{mathptmx}                      % Times (APA 7 admite Times New Roman 12 pt)
\usepackage[spanish,es-tabla]{babel}       %BABEL
\usepackage[margin=1in]{geometry}
\usepackage{amsmath,amssymb}
\usepackage{graphicx}
\usepackage{booktabs,array,multirow,tabularx}
\usepackage{setspace}
\usepackage{etoolbox}
\usepackage{enumitem}
\usepackage{xcolor}
\usepackage{csquotes}
\usepackage{tikz}
\usetikzlibrary{arrows.meta,positioning,babel}
\usepackage{algorithm}
\usepackage{algpseudocode}
\usepackage{listings}
\usepackage{caption}
\usepackage{titlesec}
\usepackage{fancyhdr}
\usepackage[hidelinks]{hyperref}

\hypersetup{pdftitle={Programación de líneas de producción: flow shop permutacional (F||Cmax)},
            pdfauthor={Camile}}

% ---------- Nombres en español -----------------------------------------------
\addto\captionsspanish{%                                                       %BABEL
  \renewcommand{\contentsname}{Índice}%                                        %BABEL
  \renewcommand{\listfigurename}{Lista de figuras}%                            %BABEL
  \renewcommand{\listtablename}{Lista de tablas}%                              %BABEL
  \renewcommand{\refname}{Referencias}%                                        %BABEL
  \renewcommand{\abstractname}{Resumen}}                                       %BABEL
\floatname{algorithm}{Algoritmo}

% ---------- Formato APA 7 ----------------------------------------------------
\setlength{\parindent}{0.5in}
\setlength{\parskip}{0pt}
\doublespacing
\setlist{itemsep=0pt,topsep=0pt,parsep=0pt}
\setcounter{secnumdepth}{0}                % APA: encabezados sin numeración
\setcounter{tocdepth}{2}

\titleformat{\section}{\normalfont\bfseries\centering}{}{0pt}{}
\titleformat{\subsection}{\normalfont\bfseries}{}{0pt}{}
\titleformat{\subsubsection}{\normalfont\bfseries\itshape}{}{0pt}{}
\titlespacing*{\section}{0pt}{0pt}{0pt}
\titlespacing*{\subsection}{0pt}{0pt}{0pt}
\titlespacing*{\subsubsection}{0pt}{0pt}{0pt}

\pagestyle{fancy}
\fancyhf{}
\fancyhead[R]{\thepage}
\renewcommand{\headrulewidth}{0pt}
\setlength{\headheight}{15pt}

\captionsetup{labelfont=bf,textfont=it,labelsep=newline,singlelinecheck=false,
              justification=raggedright,font=normalsize,skip=6pt}
\AtBeginEnvironment{table}{\singlespacing}
\AtBeginEnvironment{figure}{\singlespacing}
\AtBeginEnvironment{algorithm}{\singlespacing}

% Traducción del pseudocódigo
\algrenewcommand\algorithmicrequire{\textbf{Entrada:}}
\algrenewcommand\algorithmicensure{\textbf{Salida:}}
\algrenewcommand\algorithmicfor{\textbf{para}}
\algrenewcommand\algorithmicdo{\textbf{hacer}}
\algrenewcommand\algorithmicwhile{\textbf{mientras}}
\algrenewcommand\algorithmicif{\textbf{si}}
\algrenewcommand\algorithmicthen{\textbf{entonces}}
\algrenewcommand\algorithmicelse{\textbf{si no}}
\algrenewcommand\algorithmicreturn{\textbf{devolver}}
\algrenewcommand\algorithmicend{\textbf{fin}}

\lstset{language=Python,basicstyle=\ttfamily\footnotesize,breaklines=true,
        frame=single,numbers=left,numberstyle=\tiny,xleftmargin=1.5em,
        keywordstyle=\bfseries,commentstyle=\itshape}

% ---------- Macros propias ---------------------------------------------------
\newcommand{\relleno}[1]{\textcolor{red}{[#1]}}
\newcommand{\Cmax}{C_{\max}}
% Nota de tabla/figura (APA)
\newcommand{\nota}[1]{\par\vspace{3pt}\noindent{\singlespacing\small\textit{Nota.} #1\par}}
% Incluye la imagen si existe; si no, muestra un recuadro con instrucciones
\newcommand{\imagenopcional}[3]{%
  \IfFileExists{#1}{\includegraphics[width=#2\textwidth]{#1}}{%
    \fbox{\parbox{0.92\textwidth}{\singlespacing\small\textbf{[FIGURA PENDIENTE]} #3}}}}

\begin{document}

% =============================================================================
% PORTADA
% =============================================================================
\thispagestyle{fancy}
\begin{center}
\vspace*{1.2in}
{\bfseries Programación de líneas de producción:\\
Flow shop permutacional ($F\,\|\,\Cmax$)}\\[2\baselineskip]
Capatinta Almanza Camilo Bladimir Alexander && TERCEROS\\[0.5\baselineskip]
Escuela Profesional de Ingeniería de Sistemas\\
Facultad de Ingeniería de Producción y Servicios\\
Universidad Nacional de San Agustín de Arequipa\\[0.5\baselineskip]
\relleno{Nombre y código del curso}\\
\relleno{Docente del curso}\\
18 de noviembre de 2026
\end{center}
\newpage

% =============================================================================
% RESUMEN
% =============================================================================
\begin{center}\textbf{Resumen}\end{center}
\noindent
El flow shop permutacional con minimización del tiempo total de finalización
(\textit{makespan}), denotado $Fm\,|\,prmu\,|\,\Cmax$, es uno de los problemas
clásicos de la programación de la producción en líneas donde todos los trabajos
visitan las mismas máquinas en el mismo orden. Este informe revisa su
formulación matemática, su complejidad computacional y los principales métodos de
solución: la regla de Johnson (óptima para dos máquinas), las heurísticas
constructivas de Palmer, Campbell--Dudek--Smith (CDS) y Nawaz--Enscore--Ham
(NEH), y las metaheurísticas basadas en búsqueda tabú, recocido simulado y
\textit{iterated greedy}. Se presenta además un ejemplo numérico de cinco
trabajos y cuatro máquinas, verificado mediante enumeración exhaustiva, y se
describen los \textit{benchmarks} de Taillard y las métricas habituales de
evaluación. La literatura revisada muestra que el problema es NP-difícil para tres
o más máquinas, que NEH continúa siendo la referencia constructiva y que los
algoritmos de tipo \textit{iterated greedy} ofrecen una relación muy favorable
entre simplicidad y calidad de solución.

\noindent\textit{Palabras clave:} programación de la producción, flow shop
permutacional, makespan, regla de Johnson, heurística NEH, \textit{iterated greedy}.

\bigskip
\begin{center}\textbf{Abstract}\end{center}
\noindent
The permutation flow shop with makespan minimization, denoted
$Fm\,|\,prmu\,|\,\Cmax$, is a classical production scheduling problem in which every
job visits the same machines in the same order. This report reviews its
mathematical formulation, computational complexity and main solution methods:
Johnson's rule (optimal for two machines), the constructive heuristics of Palmer,
Campbell--Dudek--Smith (CDS) and Nawaz--Enscore--Ham (NEH), and metaheuristics based
on tabu search, simulated annealing and iterated greedy. A worked example with five
jobs and four machines, verified by exhaustive enumeration, is included, together
with Taillard's benchmarks and the usual evaluation metrics. The reviewed literature
shows that the problem is NP-hard for three or more machines, that NEH remains the
constructive reference, and that iterated greedy algorithms offer an excellent
trade-off between simplicity and solution quality.

\noindent\textit{Keywords:} production scheduling, permutation flow shop, makespan,
Johnson's rule, NEH heuristic, iterated greedy.
\newpage

% =============================================================================
% ÍNDICES
% =============================================================================
{\singlespacing
\tableofcontents
\listoffigures
\listoftables}
\newpage

% =============================================================================
% INTRODUCCIÓN
% =============================================================================
\section{Programación de líneas de producción: flow shop permutacional ($F\,\|\,\Cmax$)}

\subsection{Introducción}

La programación de la producción consiste en asignar trabajos a recursos
(máquinas, estaciones o puestos) a lo largo del tiempo, de modo que se optimice
algún criterio de desempeño como la duración total, el tiempo de flujo o el
cumplimiento de fechas de entrega (Pinedo, 2016). En las líneas de producción,
donde los productos recorren una serie de estaciones dispuestas en el mismo
orden, el modelo natural es el \textit{flow shop}, también llamado taller de flujo
o de flujo continuo. Hejazi y Saghafian (2005) lo describen como uno de los
problemas más estudiados de la programación y recopilan las contribuciones
desde el trabajo pionero de Johnson (1954) hasta las metaheurísticas de comienzos
de la década de 2000.

El interés por este problema es tanto práctico como teórico. Desde el punto de
vista práctico, muchas instalaciones reales se parecen a un flow shop, de modo que
una buena secuencia de trabajos reduce el tiempo total de la jornada, el trabajo
en proceso y los tiempos ociosos de las máquinas. Desde el punto de vista teórico,
el problema combina una definición muy simple con una dificultad computacional
considerable: para tres o más máquinas no se conoce ningún algoritmo polinomial
que garantice el óptimo (Garey et al., 1976). Una revisión sistemática reciente
sobre flow shops permutacionales encontró que alrededor del 56\,\% de los estudios
analizados optimizan el \textit{makespan}, frente a un 18\,\% de estudios
multiobjetivo y un 15\,\% centrados en el tiempo total de flujo (Koç y Güner, 2023).

El presente informe tiene como objetivo describir y analizar el problema de flow
shop permutacional con minimización del makespan, $F\,\|\,\Cmax$ en la notación de
Taillard (1993) y $Fm\,|\,prmu\,|\,\Cmax$ en la notación de tres campos. Para ello se
presentan: (a) el marco conceptual y los supuestos; (b) la formulación matemática;
(c) su complejidad computacional; (d) los métodos de solución exactos, heurísticos
y metaheurísticos; (e) un ejemplo numérico reproducible; y (f) los \textit{benchmarks}
y métricas de evaluación empleados en la literatura.

% =============================================================================
\subsection{Marco conceptual}

\subsubsection{Líneas de producción y flow shop}

En un flow shop, un conjunto de $n$ trabajos debe procesarse en $m$ máquinas
dispuestas en serie. Todos los trabajos siguen la misma ruta tecnológica
$M_1 \rightarrow M_2 \rightarrow \cdots \rightarrow M_m$, y cada operación del trabajo
$j$ en la máquina $i$ tiene una duración $p_{ji}$ (ver Figura~\ref{fig:esquema}).
Es un caso particular del \textit{job shop}, en el que las rutas pueden diferir
entre trabajos (Wikipedia, s. f.; Pinedo, 2016).

\begin{figure}[htbp]
\caption{Esquema de un flow shop con $m$ máquinas en serie}\label{fig:esquema}
\centering
% Figura 1: dibujada en TikZ (no requiere archivo externo).
\begin{tikzpicture}[>=Stealth,
  maq/.style={draw,rounded corners=2pt,minimum width=1.3cm,minimum height=0.9cm,fill=gray!15},
  lab/.style={font=\scriptsize}]
  \node (in) {Trabajos};
  \node[maq,right=1.0cm of in] (m1) {$M_1$};
  \node[maq,right=1.5cm of m1] (m2) {$M_2$};
  \node[right=1.0cm of m2] (dots) {$\cdots$};
  \node[maq,right=1.0cm of dots] (mm) {$M_m$};
  \node[right=1.0cm of mm] (out) {Terminados};
  \draw[->] (in) -- (m1);
  \draw[->] (m1) -- node[above,lab]{cola} (m2);
  \draw[->] (m2) -- (dots);
  \draw[->] (dots) -- (mm);
  \draw[->] (mm) -- (out);
\end{tikzpicture}
\nota{Elaboración propia. Cada trabajo visita las máquinas en el orden
$M_1,\ldots,M_m$; entre máquinas existe espacio de espera (\textit{buffer}) ilimitado.}
\end{figure}

\begin{figure}[htbp]
\caption{Línea de producción real como ejemplo de flow shop}\label{fig:foto}
\centering
\imagenopcional{fig/linea_produccion.jpg}{0.7}{%
Descargar una fotografía de una línea de producción/ensamblaje desde Wikimedia
Commons (por ejemplo, en la categoría \url{https://commons.wikimedia.org/wiki/Category:Assembly_lines}),
elegir una con licencia libre (CC BY, CC BY-SA o dominio público), guardarla como
\texttt{fig/linea\_produccion.jpg} y completar la nota con autor, año y licencia.
(Categoría sugerida; verifique que el enlace y la licencia estén vigentes.)}
\nota{Fotografía de \relleno{Autor}, \relleno{año}, Wikimedia Commons, \relleno{licencia}.
\relleno{URL del archivo}.}
\end{figure}

\subsubsection{Notación de tres campos}

Los problemas de programación se clasifican con la notación $\alpha\,|\,\beta\,|\,\gamma$
de Graham et al. (1979), donde $\alpha$ describe el entorno de máquinas, $\beta$ las
restricciones y $\gamma$ el criterio. En este informe, $\alpha = Fm$ (flow shop con
$m$ máquinas), $\beta = prmu$ (secuencia permutacional, es decir, el mismo orden de
trabajos en todas las máquinas) y $\gamma = \Cmax$ (tiempo máximo de finalización).
Taillard (1993) abrevia el problema como $F\,\|\,\Cmax$, y Fernandez-Viagas et al.
(2017) lo denotan $Fm\,|\,prmu\,|\,\Cmax$. La Tabla~\ref{tab:notacion} resume los símbolos
usados en el informe.

\begin{table}[htbp]
\caption{Notación utilizada}\label{tab:notacion}
\centering
\begin{tabularx}{\textwidth}{@{}l X@{}}
\toprule
Símbolo & Significado \\
\midrule
$n$, $m$ & Número de trabajos y de máquinas \\
$j = 1,\ldots,n$;\; $i = 1,\ldots,m$ & Índices de trabajo y de máquina \\
$p_{ji}$ & Tiempo de procesamiento del trabajo $j$ en la máquina $i$ \\
$\pi = (\pi(1),\ldots,\pi(n))$ & Permutación (secuencia) de los trabajos \\
$C_{\pi(k),i}$ & Tiempo de finalización del $k$-ésimo trabajo de $\pi$ en la máquina $i$ \\
$\Cmax$ & \textit{Makespan}: $C_{\pi(n),m}$ \\
$T_j$ & Tiempo total de procesamiento del trabajo $j$: $\sum_{i=1}^m p_{ji}$ \\
\bottomrule
\end{tabularx}
\nota{Elaboración propia.}
\end{table}

\subsubsection{Supuestos del problema}

El problema clásico $Fm\,|\,prmu\,|\,\Cmax$ se plantea bajo los siguientes supuestos:
\begin{itemize}
  \item Los $n$ trabajos están disponibles en el instante cero y todas las máquinas
        están libres.
  \item Cada trabajo visita las máquinas $M_1,\ldots,M_m$ en ese orden y no puede
        saltarse ninguna.
  \item Los tiempos $p_{ji}\ge 0$ son deterministas, conocidos e independientes de la
        secuencia (no hay tiempos de preparación dependientes de la secuencia).
  \item Una máquina procesa un solo trabajo a la vez y un trabajo es procesado por una
        sola máquina a la vez.
  \item No se permiten interrupciones (no hay \textit{preemption}).
  \item El espacio de espera entre máquinas es ilimitado.
  \item La secuencia de trabajos es la misma en todas las máquinas (restricción
        permutacional).
\end{itemize}
Variantes con \textit{no-wait}, bloqueo o tiempos de preparación se estudian por
separado (Hejazi y Saghafian, 2005; Ruiz y Stützle, 2008), y las versiones flexibles,
híbridas o distribuidas quedan fuera del alcance de este informe (Koç y Güner, 2023).

% =============================================================================
\subsection{Formulación matemática}

\subsubsection{Cálculo del makespan}

Dada una permutación $\pi$, los tiempos de finalización se calculan con la
recurrencia
\begin{equation}\label{eq:rec}
C_{\pi(k),i} = \max\bigl\{C_{\pi(k-1),i},\; C_{\pi(k),i-1}\bigr\} + p_{\pi(k),i},
\end{equation}
con $C_{\pi(0),i}=0$ y $C_{\pi(k),0}=0$. El makespan es $\Cmax(\pi)=C_{\pi(n),m}$. La
primera cantidad del máximo refleja que la máquina $i$ debe haber terminado el trabajo
anterior y la segunda, que el trabajo debe haber terminado en la máquina previa. Evaluar
una secuencia completa cuesta $O(nm)$ operaciones.

Equivalentemente, $\Cmax(\pi)$ es la longitud del camino más largo (camino crítico)
en un grafo de rejilla de $n\times m$ nodos, donde el nodo $(k,i)$ tiene peso
$p_{\pi(k),i}$ y los arcos van de $(k,i)$ hacia $(k+1,i)$ y hacia $(k,i+1)$
(ver Figura~\ref{fig:grafo}).

\begin{figure}[htbp]
\caption{Makespan como camino más largo en un grafo de rejilla}\label{fig:grafo}
\centering
% Figura 4: dibujada en TikZ (no requiere archivo externo).
\begin{tikzpicture}[x=2.4cm,y=-1.3cm,
  nodo/.style={draw,circle,minimum size=1.1cm,inner sep=0pt,font=\footnotesize},
  arco/.style={-{Stealth[length=2mm]},gray},
  crit/.style={-{Stealth[length=2.4mm]},line width=1.6pt,black}]
  \foreach \k in {1,2,3}{
    \node[font=\footnotesize] at (\k,0.1) {$k=\k$};
    \foreach \i in {1,2,3}{
      \node[nodo] (n\k\i) at (\k,\i) {$p_{\pi(\k),\i}$};
    }
  }
  \foreach \i in {1,2,3}{ \node[font=\footnotesize,anchor=east] at (0.55,\i) {$M_\i$}; }
  \foreach \k/\kn in {1/2,2/3}{\foreach \i in {1,2,3}{\draw[arco] (n\k\i) -- (n\kn\i);}}
  \foreach \k in {1,2,3}{\foreach \i/\ii in {1/2,2/3}{\draw[arco] (n\k\i) -- (n\k\ii);}}
  \draw[crit] (n11) -- (n21);
  \draw[crit] (n21) -- (n22);
  \draw[crit] (n22) -- (n23);
  \draw[crit] (n23) -- (n33);
\end{tikzpicture}
\nota{Elaboración propia. Ejemplo con $n=3$ y $m=3$; el trazo grueso indica un camino
crítico posible, cuya suma de pesos es $\Cmax$.}
\end{figure}

\subsubsection{Modelo de programación lineal entera mixta}

Una formulación posicional estándar emplea variables binarias $x_{jk}=1$ si el trabajo
$j$ ocupa la posición $k$ de la secuencia, y variables continuas $C_{k,i}$ con el tiempo
de finalización del trabajo situado en la posición $k$ en la máquina $i$:
\begin{align}
\min\;& \Cmax \label{eq:obj}\\
\text{s.a.}\;& \sum_{k=1}^{n} x_{jk} = 1, && j=1,\ldots,n, \label{eq:asig1}\\
& \sum_{j=1}^{n} x_{jk} = 1, && k=1,\ldots,n, \label{eq:asig2}\\
& C_{1,1} \ge \sum_{j=1}^{n} p_{j1}\,x_{j1}, \label{eq:ini}\\
& C_{k,i} \ge C_{k,i-1} + \sum_{j=1}^{n} p_{ji}\,x_{jk}, && k=1,\ldots,n;\; i=2,\ldots,m, \label{eq:maq}\\
& C_{k,i} \ge C_{k-1,i} + \sum_{j=1}^{n} p_{ji}\,x_{jk}, && k=2,\ldots,n;\; i=1,\ldots,m, \label{eq:pos}\\
& \Cmax \ge C_{n,m}, \qquad x_{jk}\in\{0,1\},\; C_{k,i}\ge 0. \label{eq:dom}
\end{align}
Las restricciones \eqref{eq:asig1}--\eqref{eq:asig2} garantizan que la secuencia sea una
permutación, y \eqref{eq:maq}--\eqref{eq:pos} reproducen la recurrencia \eqref{eq:rec}. La
programación lineal entera mixta es el enfoque de modelado más frecuente en la
literatura reciente, con un 67\,\% de los artículos que especifican un modelo, y
CPLEX fue el único \textit{solver} reportado en esa revisión (Koç y Güner, 2023).
Sin embargo, estos modelos solo resuelven instancias pequeñas o medianas, lo que
motiva el uso de heurísticas.

\subsubsection{Espacio de soluciones}

En un flow shop general, cada máquina podría procesar los trabajos en un orden distinto,
por lo que existen $(n!)^m$ programas posibles; restringirse a permutaciones reduce el
espacio a $n!$ secuencias (Hejazi y Saghafian, 2005). Para el criterio $\Cmax$ se sabe
que existe siempre una solución óptima con la misma secuencia en las dos primeras
máquinas y en las dos últimas, de modo que para $m\le 3$ las permutaciones son
suficientes (Pinedo, 2016). Para $m>3$ el óptimo permutacional puede ser peor que el
óptimo general, pero la mayor parte de la investigación se concentra en el caso
permutacional por su menor tamaño y su relevancia práctica (Hejazi y Saghafian, 2005).

Aun así, $n!$ crece muy rápido: con $n=20$ trabajos hay aproximadamente
$2{,}43\times 10^{18}$ secuencias, de modo que la enumeración completa es inviable
(ver Figura~\ref{fig:espacio}).

\begin{figure}[htbp]
\caption{Tamaño del espacio de búsqueda según el número de trabajos}\label{fig:espacio}
\centering
\imagenopcional{fig/espacio_busqueda.png}{0.75}{%
Generar con el script \texttt{generar\_figuras.py} (Python 3 + matplotlib):
\texttt{pip install matplotlib \&\& python3 generar\_figuras.py}. Crea
\texttt{fig/espacio\_busqueda.png}.}
\nota{Elaboración propia. Escala logarítmica; $m$ es el número de máquinas del caso general.}
\end{figure}

% =============================================================================
\subsection{Complejidad computacional}

Garey et al. (1976) demostraron que determinar un programa de longitud mínima en un flow
shop de $m$ máquinas es NP-completo para $m\ge 3$, mientras que para $m=2$ existe un
algoritmo eficiente; sus resultados son fuertes, es decir, valen aunque el tamaño del
problema se mida por el número de tareas o por la suma de sus duraciones. Gonzalez y Sahni
(1978) mostraron que obtener programas de finalización mínima en flow shops y job shops
es NP-completo tanto con interrupciones como sin ellas, y obtuvieron cotas de desempeño
para varias heurísticas. En consecuencia, $F2\,|\,prmu\,|\,\Cmax$ se resuelve en tiempo
$O(n\log n)$ mediante la regla de Johnson (1954), pero para $m\ge 3$ no se conoce un
algoritmo que garantice el óptimo en tiempo polinomial, y la investigación se ha
orientado hacia heurísticas y metaheurísticas (Framiñán et al., 2004).

% =============================================================================
\subsection{Métodos de solución}

\subsubsection{Método exacto para dos máquinas: regla de Johnson}

Johnson (1954) demostró que, para dos máquinas, existe una secuencia óptima en la que el
trabajo $j$ precede al trabajo $k$ si $\min\{p_{j1},p_{k2}\} \le \min\{p_{k1},p_{j2}\}$.
El Algoritmo~\ref{alg:johnson} construye esa secuencia en $O(n\log n)$: los trabajos
cuya primera operación es la más corta van primero, en orden creciente de $p_{j1}$, y los
demás van al final, en orden decreciente de $p_{j2}$.

\begin{algorithm}[htbp]
\caption{Regla de Johnson para $F2\,|\,prmu\,|\,\Cmax$}\label{alg:johnson}
\begin{algorithmic}[1]
\Require Tiempos $p_{j1}$ y $p_{j2}$ de los $n$ trabajos
\State $A \gets \{j : p_{j1} \le p_{j2}\}$;\quad $B \gets \{j : p_{j1} > p_{j2}\}$
\State Ordenar $A$ de forma creciente según $p_{j1}$
\State Ordenar $B$ de forma decreciente según $p_{j2}$
\State \Return $\pi = (A, B)$
\end{algorithmic}
\end{algorithm}

Johnson (1954) también estudió el caso de tres máquinas: si $\min_j p_{j1}\ge\max_j p_{j2}$
o $\min_j p_{j3}\ge\max_j p_{j2}$, la secuencia óptima se obtiene aplicando la regla a los
tiempos agregados $(p_{j1}+p_{j2},\;p_{j2}+p_{j3})$.

\subsubsection{Cotas inferiores}

Una cota inferior sencilla se obtiene analizando cada máquina $k$ por separado: ningún
programa puede terminar antes de que la máquina $k$ procese todos los trabajos, precedida
por el menor tiempo de llegada posible y seguida por la menor cola posible,
\begin{equation}\label{eq:lb}
LB = \max_{k=1,\ldots,m}\Bigl\{\min_j \sum_{i<k} p_{ji} \;+\; \sum_{j=1}^{n} p_{jk} \;+\; \min_j \sum_{i>k} p_{ji}\Bigr\}.
\end{equation}
Estas cotas son la base de los algoritmos de ramificación y acotamiento, que según Koç y
Güner (2023) son la técnica exacta dominante en la literatura permutacional y se emplean
también para medir la calidad de las heurísticas.

\subsubsection{Heurísticas constructivas}

Framiñán et al. (2004) revisaron y clasificaron las heurísticas para el flow shop
permutacional con objetivo $\Cmax$ y propusieron un marco común en fases (generación de
un índice de prioridad, construcción de la solución y mejora), de modo que sea posible
combinarlas y extenderlas. Las tres más conocidas son las siguientes.

\paragraph{Índice de pendiente de Palmer.} Palmer (1965) asigna a cada trabajo el índice
$s_j=\sum_{i=1}^{m}(2i-m-1)\,p_{ji}$, que es mayor para trabajos con tiempos crecientes a
lo largo de la línea, y ordena los trabajos por $s_j$ no creciente. Su costo es
$O(nm+n\log n)$.

\paragraph{CDS.} Campbell et al. (1970) generalizan la regla de Johnson: para
$k=1,\ldots,m-1$ construyen un problema de dos máquinas con tiempos
$p'_{j1}=\sum_{i=1}^{k}p_{ji}$ y $p'_{j2}=\sum_{i=m-k+1}^{m}p_{ji}$, lo resuelven con la regla
de Johnson y conservan la mejor de las $m-1$ secuencias según su makespan real.

\paragraph{NEH.} Nawaz et al. (1983) parten de la idea de que los trabajos con mayor tiempo
total de procesamiento deben recibir mayor prioridad; el algoritmo se resume en el
Algoritmo~\ref{alg:neh}. Los autores informaron que el método produce muy buenas secuencias
frente a otras heurísticas, con especial buen desempeño en problemas grandes (Nawaz et al.,
1983). Taillard (1990) mostró que la evaluación de las $k$ posiciones de inserción puede
hacerse en $O(nm)$ en total mediante cabezas $e_{l,i}$, colas $q_{l,i}$ y
$f_{l,i}=\max\{f_{l,i-1},e_{l-1,i}\}+p_{j,i}$, de modo que
$\Cmax$ al insertar $j$ en la posición $l$ es $\max_i\{f_{l,i}+q_{l,i}\}$; así, NEH pasa de
$O(n^3m)$ a $O(n^2m)$. NEH se considera todavía hoy la heurística de referencia del problema
(Fernandez-Viagas et al., 2017; Kalczynski y Kamburowski, 2007), y fue la heurística más
usada en los estudios revisados por Koç y Güner (2023). Framiñán et al. (2003) estudiaron
órdenes iniciales alternativos para NEH, y Ruiz y Maroto (2005) la incluyeron como
referencia en su evaluación comparativa de heurísticas permutacionales.

\begin{algorithm}[htbp]
\caption{Heurística NEH}\label{alg:neh}
\begin{algorithmic}[1]
\Require Matriz de tiempos $p_{ji}$
\State $T_j \gets \sum_{i=1}^m p_{ji}$ para todo $j$
\State Ordenar los trabajos de forma no creciente según $T_j$: $(\sigma_1,\ldots,\sigma_n)$
\State $\pi \gets (\sigma_1)$
\For{$k = 2,\ldots,n$}
  \State Insertar $\sigma_k$ en cada una de las $k$ posiciones de $\pi$ y calcular $\Cmax$
  \State $\pi \gets$ la secuencia parcial con menor $\Cmax$ (empates: primera posición)
\EndFor
\State \Return $\pi$
\end{algorithmic}
\end{algorithm}

\subsubsection{Metaheurísticas}

Las metaheurísticas parten de una solución inicial (típicamente la de NEH) y la mejoran
mediante búsqueda local guiada. Osman y Potts (1989) aplicaron recocido simulado al flow
shop permutacional, y Taillard (1990) y Nowicki y Smutnicki (1996) propusieron algoritmos
de búsqueda tabú; este último se presenta como un algoritmo rápido para el problema. Los
algoritmos genéticos y los híbridos con heurísticas constructivas también son frecuentes:
Bari y Karande (2023) propusieron un algoritmo genético cuya población inicial se siembra
con secuencias de Palmer, CDS, NEH y PSO, y lo aplicaron a datos de una siderúrgica; en
su caso de estudio redujeron el makespan en alrededor de un 16\,\% respecto a la regla
de atender por orden de llegada de la empresa, aunque igualaron el resultado de NEH en ese
caso. En la revisión de Koç y Güner (2023), las metaheurísticas fueron el enfoque de
solución más común (68\,\% de los artículos), seguidas de las heurísticas (21\,\%) y los
métodos exactos (6\,\%).

Entre todas ellas destaca el algoritmo \textit{iterated greedy} (IG) de Ruiz y Stützle
(2007), que alterna una fase de destrucción, en la que se eliminan $d$ trabajos de la
solución actual, con una fase de construcción, en la que se reinsertan usando el
procedimiento de inserción de NEH, y que opcionalmente aplica una búsqueda local posterior.
Los autores reportan que es muy simple de implementar y altamente eficaz frente a los
métodos del estado del arte de su momento. El Algoritmo~\ref{alg:ig} resume su esquema;
en el trabajo original se emplea una temperatura constante
$T=T_p\cdot\sum_{j}\sum_{i}p_{ji}\,/\,(10\,nm)$ con $d=4$ y $T_p=0.4$ como valores
recomendados. Posteriormente, Dubois-Lacoste et al. (2017) propusieron una variante de
IG que optimiza soluciones parciales para el mismo problema.

\begin{algorithm}[htbp]
\caption{Esquema del algoritmo \textit{iterated greedy} (IG)}\label{alg:ig}
\begin{algorithmic}[1]
\Require Matriz $p_{ji}$, tamaño de destrucción $d$, temperatura $T$
\State $\pi \gets \text{NEH}(p)$;\quad $\pi^{*} \gets \pi$
\While{no se cumpla el criterio de parada}
  \State $\pi' \gets \pi$;\quad eliminar $d$ trabajos al azar de $\pi'$ (lista $D$)
  \For{cada trabajo $j \in D$}
    \State insertar $j$ en la posición de $\pi'$ que minimiza $\Cmax$
  \EndFor
  \State (opcional) aplicar búsqueda local por inserción a $\pi'$
  \If{$\Cmax(\pi') < \Cmax(\pi)$}
    \State $\pi \gets \pi'$;\quad \If{$\Cmax(\pi') < \Cmax(\pi^{*})$} $\pi^{*}\gets\pi'$ \EndIf
  \ElsIf{$\text{aleatorio}(0,1) \le \exp\bigl(-(\Cmax(\pi')-\Cmax(\pi))/T\bigr)$}
    \State $\pi \gets \pi'$
  \EndIf
\EndWhile
\State \Return $\pi^{*}$
\end{algorithmic}
\end{algorithm}

La Tabla~\ref{tab:metodos} resume los métodos presentados.

\begin{table}[htbp]
\caption{Resumen de los principales métodos para $Fm\,|\,prmu\,|\,\Cmax$}\label{tab:metodos}
\centering
\small
\begin{tabularx}{\textwidth}{@{}l l l X@{}}
\toprule
Método & Año & Tipo & Característica principal \\
\midrule
Regla de Johnson & 1954 & Exacto ($m=2$) & $O(n\log n)$; óptimo para dos máquinas \\
Palmer & 1965 & Constructivo & Índice de pendiente; $O(nm+n\log n)$ \\
CDS & 1970 & Constructivo & $m-1$ problemas de dos máquinas resueltos con Johnson \\
NEH & 1983 & Constructivo & Inserción por tiempo total decreciente; $O(n^2m)$ con aceleración \\
Búsqueda tabú & 1990--1996 & Metaheurística & Vecindario de inserción con memoria tabú \\
Recocido simulado & 1989 & Metaheurística & Aceptación probabilística de soluciones peores \\
\textit{Iterated greedy} & 2007 & Metaheurística & Destrucción--construcción con reinserción tipo NEH \\
\bottomrule
\end{tabularx}
\nota{Elaboración propia a partir de Johnson (1954), Palmer (1965), Campbell et al. (1970),
Nawaz et al. (1983), Taillard (1990), Osman y Potts (1989), Nowicki y Smutnicki (1996) y
Ruiz y Stützle (2007).}
\end{table}

% =============================================================================
\subsection{Ejemplos numéricos}

Los resultados de esta sección se calcularon con el código del Apéndice y se verificaron
por enumeración exhaustiva de todas las permutaciones.

\subsubsection{Dos máquinas: regla de Johnson}

Considérense cinco trabajos (A, B, C, D, E) en dos máquinas con los tiempos de la
Tabla~\ref{tab:johnson}. El conjunto $A=\{\text{A},\text{C},\text{D}\}$ (con
$p_{j1}\le p_{j2}$) se ordena por $p_{j1}$ creciente: C, A, D; el conjunto
$B=\{\text{B},\text{E}\}$ se ordena por $p_{j2}$ decreciente: E, B. La secuencia resultante
es C--A--D--E--B, con $\Cmax=24$ (Figura~\ref{fig:johnson}). Como la suma de $p_{j1}$ es
22 y el menor $p_{j2}$ es 2, la cota inferior es 24, por lo que la secuencia es óptima.

\begin{table}[htbp]
\caption{Tiempos de procesamiento del ejemplo de dos máquinas}\label{tab:johnson}
\centering
\begin{tabular}{@{}lccccc@{}}
\toprule
Trabajo & A & B & C & D & E \\
\midrule
$p_{j1}$ & 3 & 5 & 1 & 6 & 7 \\
$p_{j2}$ & 6 & 2 & 2 & 6 & 5 \\
\bottomrule
\end{tabular}
\nota{Elaboración propia.}
\end{table}

\begin{figure}[htbp]
\caption{Diagrama de Gantt de la secuencia de Johnson}\label{fig:johnson}
\centering
\imagenopcional{fig/gantt_johnson.png}{0.85}{%
Generar con \texttt{python3 generar\_figuras.py} (Python 3 + matplotlib); crea
\texttt{fig/gantt\_johnson.png}.}
\nota{Elaboración propia. La línea discontinua marca $\Cmax=24$.}
\end{figure}

\subsubsection{Cuatro máquinas: Palmer, CDS y NEH}

Se consideran cinco trabajos en cuatro máquinas con los tiempos de la
Tabla~\ref{tab:datos}.

\begin{table}[htbp]
\caption{Tiempos de procesamiento del ejemplo de cuatro máquinas}\label{tab:datos}
\centering
\begin{tabular}{@{}lccccc@{}}
\toprule
Trabajo & $M_1$ & $M_2$ & $M_3$ & $M_4$ & $T_j$ \\
\midrule
J1 & 5 & 9 & 8 & 10 & 32 \\
J2 & 9 & 3 & 10 & 1 & 23 \\
J3 & 9 & 4 & 5 & 8 & 26 \\
J4 & 4 & 8 & 8 & 7 & 27 \\
J5 & 1 & 5 & 6 & 8 & 20 \\
\bottomrule
\end{tabular}
\nota{Elaboración propia.}
\end{table}

\paragraph{NEH.} El orden inicial por $T_j$ no creciente es J1, J4, J3, J2, J5. La
Tabla~\ref{tab:neh} muestra el makespan de cada inserción; en el primer paso hay empate
y se elige la primera posición. La secuencia final es J5--J4--J3--J1--J2 con $\Cmax=48$.

\begin{table}[htbp]
\caption{Evaluación de posiciones de inserción en NEH}\label{tab:neh}
\centering
\begin{tabular}{@{}clccccl@{}}
\toprule
Paso & Trabajo & Pos. 1 & Pos. 2 & Pos. 3 & Pos. 4 & Pos. 5 \\
\midrule
1 & J4 & \textbf{39} & 39 & --- & --- & --- \\
2 & J3 & 48 & \textbf{45} & 47 & --- & --- \\
3 & J2 & 55 & 54 & 54 & \textbf{46} & --- \\
4 & J5 & \textbf{48} & 54 & 54 & 54 & 59 \\
\bottomrule
\end{tabular}
\nota{Elaboración propia. Valores de $\Cmax$ de la secuencia parcial; en negrita, la posición
elegida. Secuencias parciales elegidas: J4--J1, J4--J3--J1, J4--J3--J1--J2 y
J5--J4--J3--J1--J2.}
\end{table}

La Tabla~\ref{tab:comp} muestra los tiempos de finalización de la secuencia NEH calculados
con la recurrencia \eqref{eq:rec}.

\begin{table}[htbp]
\caption{Tiempos de finalización $C_{\pi(k),i}$ de la secuencia NEH}\label{tab:comp}
\centering
\begin{tabular}{@{}lcccc@{}}
\toprule
Trabajo & $M_1$ & $M_2$ & $M_3$ & $M_4$ \\
\midrule
J5 & 1 & 6 & 12 & 20 \\
J4 & 5 & 14 & 22 & 29 \\
J3 & 14 & 18 & 27 & 37 \\
J1 & 19 & 28 & 36 & 47 \\
J2 & 28 & 31 & 46 & \textbf{48} \\
\bottomrule
\end{tabular}
\nota{Elaboración propia.}
\end{table}

\paragraph{Palmer y CDS.} Con $m=4$, los coeficientes de Palmer son $-3,-1,1,3$ y los
índices resultan $s=(14,-17,-2,9,22)$ para J1,\ldots,J5, lo que da la secuencia
J5--J1--J4--J3--J2 con $\Cmax=49$. En CDS, los tres problemas auxiliares ($k=1,2,3$)
producen makespans de 50, 48 y 50; la mejor secuencia, J5--J4--J3--J1--J2, tiene
$\Cmax=48$.

\paragraph{Comparación.} La enumeración de las $5!=120$ permutaciones confirma que 48 es
el óptimo. La cota inferior de la ecuación \eqref{eq:lb} es 46. La Tabla~\ref{tab:res}
resume los resultados y la Figura~\ref{fig:gantt} compara el Gantt del orden original con
el de NEH.

\begin{table}[htbp]
\caption{Resultados del ejemplo de cuatro máquinas}\label{tab:res}
\centering
\begin{tabular}{@{}llcc@{}}
\toprule
Método & Secuencia & $\Cmax$ & Desvío respecto al óptimo (\%) \\
\midrule
Orden original & J1--J2--J3--J4--J5 & 60 & 25.0 \\
Palmer & J5--J1--J4--J3--J2 & 49 & 2.1 \\
CDS & J5--J4--J3--J1--J2 & 48 & 0.0 \\
NEH & J5--J4--J3--J1--J2 & 48 & 0.0 \\
Óptimo (enumeración) & J5--J4--J3--J1--J2 & 48 & 0.0 \\
\bottomrule
\end{tabular}
\nota{Elaboración propia. Cota inferior: 46.}
\end{table}

\begin{figure}[htbp]
\caption{Diagramas de Gantt del orden original y de la secuencia NEH}\label{fig:gantt}
\centering
\imagenopcional{fig/gantt_comparacion.png}{0.9}{%
Generar con \texttt{python3 generar\_figuras.py} (Python 3 + matplotlib); crea
\texttt{fig/gantt\_comparacion.png}.}
\nota{Elaboración propia. Cada barra es una operación; los espacios vacíos son tiempos
ociosos de las máquinas.}
\end{figure}

% =============================================================================
\subsection{Evaluación experimental}

\subsubsection{Instancias de referencia}

Taillard (1993) propuso 260 problemas generados aleatoriamente, de tamaño comparable al de
problemas industriales, para flow shop permutacional, job shop y open shop con criterio
$\Cmax$. Para el flow shop se dispone de 120 instancias (ta001--ta120), que Beasley (1990)
distribuye a través de la OR-Library junto con el código de generación (ver
\url{https://people.brunel.ac.uk/~mastjjb/jeb/orlib/flowshopinfo.html}). Las instancias se
agrupan en doce clases de diez instancias cada una (Tabla~\ref{tab:taillard}), con tiempos
de procesamiento enteros generados uniformemente en el rango 1--99.

\begin{table}[htbp]
\caption{Clases de instancias de Taillard para el flow shop permutacional}\label{tab:taillard}
\centering
\begin{tabular}{@{}lcccccccccccc@{}}
\toprule
$n$ & 20 & 20 & 20 & 50 & 50 & 50 & 100 & 100 & 100 & 200 & 200 & 500 \\
$m$ & 5 & 10 & 20 & 5 & 10 & 20 & 5 & 10 & 20 & 10 & 20 & 20 \\
\bottomrule
\end{tabular}
\nota{Elaboración propia a partir de Taillard (1993). Cada clase contiene diez instancias,
para un total de 120.}
\end{table}

\subsubsection{Métricas}

La calidad de un algoritmo se mide habitualmente con la desviación porcentual relativa
\begin{equation}\label{eq:rpd}
RPD = \frac{\Cmax^{\text{alg}} - \Cmax^{\text{ref}}}{\Cmax^{\text{ref}}}\times 100,
\end{equation}
donde $\Cmax^{\text{ref}}$ es la mejor solución conocida o una cota inferior, promediada
sobre las instancias de cada clase (ARPD), junto con el tiempo de CPU. Fernandez-Viagas et al.
(2017) realizaron una evaluación computacional de los métodos aproximados publicados para este
problema, y Ruiz y Stützle (2007) validaron su algoritmo IG con análisis estadísticos sobre
las instancias de Taillard. En cuanto a herramientas, los estudios revisados por Koç y Güner
(2023) implementan sus algoritmos principalmente en C++ (53\,\%) y MATLAB (34\,\%).

% =============================================================================
\subsection{Discusión}

La literatura revisada permite extraer varias ideas. Primero, existe una frontera nítida
entre $m=2$, donde la regla de Johnson es óptima y eficiente, y $m\ge 3$, donde el problema es
NP-difícil (Garey et al., 1976), lo que explica la abundancia de heurísticas. Segundo,
NEH combina simplicidad y calidad: en el ejemplo de este informe alcanza el óptimo y, en la
literatura, sigue siendo la base sobre la que se construyen variantes y metaheurísticas
como IG (Fernandez-Viagas et al., 2017; Ruiz y Stützle, 2007). Tercero, el
buen desempeño en un caso de estudio no garantiza superioridad general: en el estudio de
Bari y Karande (2023), un algoritmo híbrido igualó a NEH en el caso real y solo lo superó
en una parte de las instancias aleatorias.

Entre las limitaciones del modelo estudiado se encuentran los supuestos deterministas, la
ausencia de tiempos de preparación dependientes de la secuencia y el uso de un solo
criterio. Las extensiones más relevantes incluyen tiempos de preparación (Ruiz y Stützle,
2008), restricciones de bloqueo o \textit{no-wait} (Hejazi y Saghafian, 2005) y objetivos
múltiples, que Koç y Güner (2023) identifican como una línea con vacíos por explorar. Por
último, el ejemplo numérico de este informe es pequeño y su único fin es didáctico; las
conclusiones sobre desempeño deben contrastarse con las instancias de Taillard.

% =============================================================================
\subsection{Conclusiones}

El flow shop permutacional con objetivo $\Cmax$ modela de forma directa la programación
de líneas de producción en serie. Para dos máquinas, la regla de Johnson (1954) entrega el
óptimo en $O(n\log n)$; para tres o más máquinas el problema es NP-difícil (Garey et al.,
1976), por lo que se recurre a heurísticas constructivas como Palmer, CDS y NEH, y a
metaheurísticas como búsqueda tabú, recocido simulado e \textit{iterated greedy}. En el
ejemplo de cuatro máquinas, NEH y CDS alcanzaron el óptimo de 48 unidades de tiempo, frente
a 60 del orden original, lo que ilustra el impacto de una buena secuencia. Se recomienda
implementar NEH con la aceleración de Taillard como punto de partida, extenderlo a un
algoritmo \textit{iterated greedy} y evaluarlo con las instancias de Taillard y la métrica
RPD.

% =============================================================================
% REFERENCIAS (APA 7, sangría francesa)
% =============================================================================
\newpage
\section{Referencias}
\begin{list}{}{\setlength{\leftmargin}{0.5in}\setlength{\itemindent}{-0.5in}%
  \setlength{\itemsep}{0pt}\setlength{\parsep}{0pt}\setlength{\topsep}{0pt}}

\item Bari, P., y Karande, P. (2023). Makespan minimization using hybrid heuristic and
evolutionary genetic algorithm. \textit{International Journal of Industrial Engineering \&
Production Research, 34}(2), 1--16. \url{https://doi.org/10.22068/ijiepr.34.2.8}

\item Beasley, J. E. (1990). OR-Library: Distributing test problems by electronic mail.
\textit{Journal of the Operational Research Society, 41}(11), 1069--1072.
\url{https://doi.org/10.1057/jors.1990.166}

\item Campbell, H. G., Dudek, R. A., y Smith, M. L. (1970). A heuristic algorithm for the n
job, m machine sequencing problem. \textit{Management Science, 16}(10), B630--B637.
\url{https://doi.org/10.1287/mnsc.16.10.B630}

\item Dubois-Lacoste, J., Pagnozzi, F., y Stützle, T. (2017). An iterated greedy algorithm
with optimization of partial solutions for the makespan permutation flowshop problem.
\textit{Computers \& Operations Research, 81}, 160--166.
\url{https://doi.org/10.1016/j.cor.2016.12.021}

\item Fernandez-Viagas, V., Ruiz, R., y Framinan, J. M. (2017). A new vision of approximate
methods for the permutation flowshop to minimise makespan: State-of-the-art and computational
evaluation. \textit{European Journal of Operational Research, 257}(3), 707--721.
\url{https://doi.org/10.1016/j.ejor.2016.09.055}

\item Framiñán, J. M., Gupta, J. N. D., y Leisten, R. (2004). A review and classification of
heuristics for permutation flow-shop scheduling with makespan objective. \textit{Journal of the
Operational Research Society, 55}(12), 1243--1255.
\url{https://doi.org/10.1057/palgrave.jors.2601784}

\item Framiñán, J. M., Leisten, R., y Rajendran, C. (2003). Different initial sequences for the
heuristic of Nawaz, Enscore and Ham to minimize makespan, idletime or flowtime in the static
permutation flowshop sequencing problem. \textit{International Journal of Production Research,
41}(1), 121--148. \url{https://doi.org/10.1080/00207540210161650}

\item Garey, M. R., Johnson, D. S., y Sethi, R. (1976). The complexity of flowshop and jobshop
scheduling. \textit{Mathematics of Operations Research, 1}(2), 117--129.
\url{https://doi.org/10.1287/moor.1.2.117}

\item Gonzalez, T., y Sahni, S. (1978). Flowshop and jobshop schedules: Complexity and
approximation. \textit{Operations Research, 26}(1), 36--52.
\url{https://doi.org/10.1287/opre.26.1.36}

\item Graham, R. L., Lawler, E. L., Lenstra, J. K., y Rinnooy Kan, A. H. G. (1979).
Optimization and approximation in deterministic sequencing and scheduling: A survey.
\textit{Annals of Discrete Mathematics, 5}, 287--326.
\url{https://doi.org/10.1016/S0167-5060(08)70356-X}

\item Hejazi, S. R., y Saghafian, S. (2005). Flowshop-scheduling problems with makespan
criterion: A review. \textit{International Journal of Production Research, 43}(14),
2895--2929.

\item Johnson, S. M. (1954). Optimal two- and three-stage production schedules with setup
times included. \textit{Naval Research Logistics Quarterly, 1}(1), 61--68.
\url{https://doi.org/10.1002/nav.3800010110}

\item Kalczynski, P. J., y Kamburowski, J. (2007). On the NEH heuristic for minimizing the
makespan in permutation flow shops. \textit{Omega, 35}(1), 53--60.

\item Koç, H. V., y Güner, E. (2023). A systematic literature review and an integrated
framework of the flowshop scheduling problem. \textit{Verimlilik Dergisi, 57}(3), 577--594.
\url{https://doi.org/10.51551/verimlilik.1207259}

\item Nawaz, M., Enscore, E. E., Jr., y Ham, I. (1983). A heuristic algorithm for the
m-machine, n-job flow-shop sequencing problem. \textit{Omega, 11}(1), 91--95.
\url{https://doi.org/10.1016/0305-0483(83)90088-9}

\item Nowicki, E., y Smutnicki, C. (1996). A fast tabu search algorithm for the permutation
flow-shop problem. \textit{European Journal of Operational Research, 91}(1), 160--175.
\url{https://doi.org/10.1016/0377-2217(95)00037-2}

\item Osman, I. H., y Potts, C. N. (1989). Simulated annealing for permutation flow-shop
scheduling. \textit{Omega, 17}(6), 551--557.

\item Palmer, D. S. (1965). Sequencing jobs through a multi-stage process in the minimum total
time---A quick method of obtaining a near optimum. \textit{Operational Research Quarterly,
16}(1), 101--107.

\item Pinedo, M. L. (2016). \textit{Scheduling: Theory, algorithms, and systems} (5.ª ed.).
Springer. \url{https://doi.org/10.1007/978-3-319-26580-3}

\item Ruiz, R., y Maroto, C. (2005). A comprehensive review and evaluation of permutation
flowshop heuristics. \textit{European Journal of Operational Research, 165}(2), 479--494.
\url{https://doi.org/10.1016/j.ejor.2004.04.017}

\item Ruiz, R., y Stützle, T. (2007). A simple and effective iterated greedy algorithm for the
permutation flowshop scheduling problem. \textit{European Journal of Operational Research,
177}(3), 2033--2049.

\item Ruiz, R., y Stützle, T. (2008). An iterated greedy heuristic for the sequence dependent
setup times flowshop problem with makespan and weighted tardiness objectives.
\textit{European Journal of Operational Research, 187}(3), 1143--1159.
\url{https://doi.org/10.1016/j.ejor.2006.07.029}

\item Taillard, E. (1990). Some efficient heuristic methods for the flow shop sequencing
problem. \textit{European Journal of Operational Research, 47}(1), 65--74.
\url{https://doi.org/10.1016/0377-2217(90)90090-X}

\item Taillard, E. (1993). Benchmarks for basic scheduling problems. \textit{European Journal
of Operational Research, 64}(2), 278--285.
\url{https://doi.org/10.1016/0377-2217(93)90182-M}

\item Wikipedia. (s. f.). \textit{Flow-shop scheduling}. Recuperado el 18 de noviembre de
2026, de \url{https://en.wikipedia.org/wiki/Flow-shop_scheduling}

\end{list}

% =============================================================================
% APÉNDICE
% =============================================================================
\newpage
\section{Apéndice}
\subsection{Código en Python del ejemplo numérico}

\begin{singlespace}
\begin{lstlisting}
def makespan(seq, p):
    """Cmax de una secuencia: recurrencia C[k][i]."""
    m = len(p[0])
    c = [0] * m
    for j in seq:
        c[0] += p[j][0]
        for i in range(1, m):
            c[i] = max(c[i], c[i - 1]) + p[j][i]
    return c[-1]

def johnson(a, b):
    """Regla de Johnson para dos maquinas (a = p1, b = p2)."""
    s1 = sorted([j for j in range(len(a)) if a[j] <= b[j]], key=lambda j: a[j])
    s2 = sorted([j for j in range(len(a)) if a[j] > b[j]], key=lambda j: -b[j])
    return s1 + s2

def neh(p):
    """Heuristica NEH (empates: primera posicion)."""
    orden = sorted(range(len(p)), key=lambda j: -sum(p[j]))
    seq = [orden[0]]
    for j in orden[1:]:
        mejor = None
        for pos in range(len(seq) + 1):
            cand = seq[:pos] + [j] + seq[pos:]
            c = makespan(cand, p)
            if mejor is None or c < mejor[0]:
                mejor = (c, cand)
        seq = mejor[1]
    return seq

P = [[5, 9, 8, 10], [9, 3, 10, 1], [9, 4, 5, 8], [4, 8, 8, 7], [1, 5, 6, 8]]
s = neh(P)
print([j + 1 for j in s], makespan(s, P))   # [5, 4, 3, 1, 2] 48
\end{lstlisting}
\end{singlespace}

\noindent El script completo \texttt{generar\_figuras.py} (que además implementa Palmer,
CDS, la cota inferior y la enumeración exhaustiva, y genera las figuras) acompaña a este
documento.

\end{document}